\documentclass[12pt]{article}
\usepackage{geometry}
\geometry{letterpaper, margin=1in}
\usepackage{hyperref}
\usepackage[usenames,dvipsnames]{xcolor}
\usepackage{microtype}
\usepackage{enumitem}

\definecolor{accent}{RGB}{26, 60, 110}

\begin{document}

%----------HEADER----------
\begin{center}
  {\LARGE\textbf{\textcolor{accent}{SITONG LI}}} \\[4pt]
  \small Ann Arbor, MI
  \;\textbullet\;
  \href{mailto:ruc.sitongli@gmail.com}{ruc.sitongli@gmail.com}
  \;\textbullet\;
  (952) 201-4026
\end{center}

\vspace{1em}

\noindent May 21, 2026

\vspace{1em}

\noindent Dear DECDG Hiring Committee,

\vspace{1em}

I am writing to apply for the Data Engineer position in the Development Data Group at the World Bank. I am currently completing a full-time predoctoral research fellowship at the Ross School of Business, University of Michigan, and am available to relocate to Washington, DC immediately. The World Bank's mission of putting data to work for development --- and DECDG's specific work on scaling data production, governance, and dissemination infrastructure --- is precisely the kind of high-impact, technical and mission-driven environment I want to contribute to.

My core data engineering experience is in designing and implementing end-to-end ETL pipelines that transform large volumes of unstructured source data into validated, structured, dissemination-ready statistical datasets. At Michigan Ross, I architected a scalable data production pipeline combining LLMs, fuzzy string matching, and word embeddings to extract, validate, and link records across heterogeneous administrative datasets. As a research assistant to Prof.\ William Cassidy at Washington University in St.\ Louis, I built a multi-stage ETL pipeline to ingest scanned archival documents via OCR, extract structured NLP entities using GPT-4, and load validated outputs into a version-controlled analytical dataset --- developing reusable Python functions with automated data quality checks at each stage. In a concurrent project with Prof.\ Haiwen Zhang at the University of Minnesota, I designed a reusable Python ingestion and classification pipeline for 10,000+ SEC regulatory filings from EDGAR, producing firm-level statistical indicators at scale with documented metadata schemas and Git-managed code repositories. Across all three settings, I have owned the full data production lifecycle: pipeline architecture, modular function development, quality validation, version control, and delivery of analysis-ready outputs.

My technical stack maps well to DECDG's requirements. I am proficient in Python and R for pipeline development and statistical indicator production, and have hands-on experience with Databricks and Spark for distributed data processing, Git and GitHub for version control and reproducible workflows, and AWS for cloud-based data infrastructure. I have worked extensively with large-scale structured and unstructured data --- including regulatory and administrative datasets analogous to the statistical and development indicator datasets DECDG manages --- and have experience building automated, CI/CD-oriented workflows for multi-stage data production. I am a fast learner on new platforms and institutional frameworks, and I bring the analytical and communication skills to serve as an effective liaison between technical and business stakeholders.

I am deeply motivated by the World Bank's development mission, and I am excited by the opportunity to contribute to the Data360 framework and help build the infrastructure that enables better development data to reach the people and institutions who need it. I would welcome the chance to discuss how my background fits DECDG's current data engineering priorities.

\vspace{1em}

\noindent Sincerely,\\[0.5em]
\textbf{Sitong Li}

\end{document}
